% ============================================================
% drafting.tex
% PAP expansion sections for "The Optics of Reform"
% These subsections are ready to be incorporated into main.tex
% ============================================================

% ------------------------------------------------------------------
% A. POWER ANALYSIS AND SAMPLE SIZE JUSTIFICATION
% ------------------------------------------------------------------
\subsection{Power Analysis and Sample Size}

We pre-specify a minimum detectable effect (MDE) analysis to justify our target sample of $N = 1{,}200$ recruited via Prolific.
All calculations assume a two-sided test with $\alpha = 0.05$ and power $1 - \beta = 0.80$.

\paragraph{Design and allocation.}
The experiment uses three equally sized arms (Control, Race Treatment, Gender Treatment), each assigned approximately 400 respondents.
We oversample Black respondents to constitute at least 30\% of the full sample ($n \approx 360$), ensuring adequate power for the subgroup test in Hypothesis 1.
Block randomization is conducted within strata defined by race and party identification, following the procedure detailed in the survey instrument.

\paragraph{Outcome variance assumption.}
Our primary outcomes are 7-point Likert items.
Prior survey experiments on police trust report standard deviations in the range of 1.3--1.8 scale points \citep{Stauffer2023, Peyton2022}.
We adopt a conservative assumption of $\sigma = 1.5$ for the main power calculations and conduct sensitivity analyses under $\sigma \in \{1.0, 1.5, 2.0\}$.

\paragraph{Main effect MDE.}
For a two-sample comparison between one treatment arm and the control ($n_1 = n_2 = 400$), the MDE is:
\begin{equation}
    \text{MDE} = (z_{\alpha/2} + z_{\beta})\,\sigma\sqrt{\frac{1}{n_1} + \frac{1}{n_2}}
               = (1.96 + 0.84)(1.5)\sqrt{\frac{1}{400} + \frac{1}{400}}
               \approx 0.30 \text{ scale points}.
\end{equation}
This corresponds to a standardized effect size of $d \approx 0.20$, which is below the smallest substantively meaningful effect identified in related work \citep{Stauffer2023}.

\paragraph{Subgroup MDE.}
For Hypothesis 1 (Black respondents only), the relevant subsample consists of Black respondents in the Race Treatment arm and the Control arm.
Assuming 30\% Black respondents in each arm ($n_{\text{sub}} \approx 120$ per arm), the subgroup MDE is approximately $0.54$ scale points ($d \approx 0.36$).
For Hypothesis 2 (women respondents), we assume women constitute approximately 50\% of the sample ($n_{\text{sub}} \approx 200$ per arm), yielding an MDE of approximately $0.42$ scale points ($d \approx 0.28$).

\paragraph{Pilot study.}
Prior to fielding the full study, we will conduct a pilot with $n = 100$ respondents (not included in the analysis) to: (1) assess treatment uptake and vignette comprehension, (2) estimate the empirical standard deviation of the primary outcomes, and (3) update MDE estimates using the observed $\hat{\sigma}$.
If the pilot $\hat{\sigma}$ substantially exceeds 1.5, we will increase $N$ accordingly before pre-registration closes.

\paragraph{Covariate adjustment.}
The primary OLS specification includes pre-treatment covariates (see Section~\ref{sec:empirical-strategy}).
Covariate adjustment is expected to reduce the residual variance, effectively improving power beyond the unadjusted calculations above.
We treat the unadjusted MDE as a conservative bound.

% ------------------------------------------------------------------
% B. PRE-SPECIFIED HETEROGENEOUS TREATMENT EFFECTS
% ------------------------------------------------------------------
\subsection{Pre-Specified Heterogeneous Treatment Effects}
\label{sec:hte}

Beyond the subgroup hypotheses specified in H1 and H2, we pre-specify an exploratory analysis of heterogeneous treatment effects (HTEs) using a Generalized Random Forest \citep[GRF;][]{Athey2019}.
This approach estimates conditional average treatment effects (CATEs) without requiring us to pre-specify the functional form of interactions, and naturally addresses concerns about multiple testing by constructing a single model over all covariates simultaneously.

\paragraph{Covariates of interest.}
We examine treatment effect heterogeneity along the following pre-specified dimensions:
\begin{enumerate}
    \item \textit{Race/ethnicity} (primary, Hypotheses 1 as defined above)
    \item \textit{Gender} (primary, Hypothesis 2)
    \item \textit{Party identification} (Democrat, Republican, Independent/Lean/Neither) --- partisanship may moderate receptiveness to diversity-framed institutional signals
    \item \textit{Prior police contact} --- respondents who report personal or family use-of-force experience may exhibit attenuated treatment effects due to entrenched distrust
    \item \textit{Pre-treatment diversity attitudes} --- higher baseline support for diversity in policing may produce ceiling effects or amplify treatment responses
    \item \textit{Political ideology} --- liberal respondents may be more receptive to the reform framing
\end{enumerate}

\paragraph{Estimation.}
We estimate a causal forest using the \texttt{grf} package in R \citep{tibshirani2022grf}.
For each treatment contrast (Race vs.\ Control; Gender vs.\ Control), we estimate a separate causal forest with outcomes $\{$\texttt{Trust\_General}, \texttt{Trust\_Self}, \texttt{Trust\_Hiring}$\}$.
Variance estimation uses the built-in honest splitting and out-of-bag predictions.
We summarize CATE heterogeneity using the best linear projection \citep[BLP;][]{Semenova2021}, which regresses estimated CATEs on the covariates of interest and is robust to model misspecification.
The BLP coefficients and their standard errors are reported in the main text.

\paragraph{Confirmatory vs.\ exploratory.}
HTE results are classified as \textit{exploratory}. The confirmatory hypotheses remain H1 and H2 as stated.
No multiplicity correction is applied within the HTE analysis; we report effect sizes and confidence intervals and note that these findings should be replicated.

% ------------------------------------------------------------------
% C. MECHANISM TESTS
% ------------------------------------------------------------------
\subsection{Mechanism Tests}
\label{sec:mechanisms}

We hypothesize three causal mechanisms through which the diversity frame may affect trust.
These mechanisms generate distinct, testable empirical implications, which we pre-specify below.

\paragraph{Mechanism 1 --- Identity Activation.}
\textit{Prediction:} The treatment effect on trust is mediated by an increase in perceived procedural justice.
Black respondents assigned to the Race Treatment will report higher \texttt{Procedural\_Justice} scores than controls, and this increase will partially or fully account for the trust effect.
We test this using the \texttt{mediation} package in R \citep{Tingley2014}, estimating the average causal mediation effect (ACME) and direct effect with 1,000 bootstrap draws under the sequential ignorability assumption.
We conduct the analysis separately for Black respondents (H1 mechanism) and women respondents (H2 mechanism).

\paragraph{Mechanism 2 --- Reform Intent Signal.}
\textit{Prediction:} Any signal of institutional reform --- regardless of the specific diversity frame --- increases trust through a general willingness-to-change inference.
Under this mechanism, the effect on \texttt{Police\_Resources} should be positive and of similar magnitude across both treatment arms.
We test this by comparing the treatment coefficients on \texttt{Police\_Resources} across arms using a seemingly unrelated regression or a joint Wald test.
A similar effect size across both treatments is consistent with the reform signal mechanism; a larger effect for the matched-identity treatment (Black respondents in Race arm; women in Gender arm) would be more consistent with Mechanism 1.

\paragraph{Mechanism 3 --- Cynicism.}
\textit{Prediction:} Among respondents with high prior distrust or who perceive reform signals as symbolic, the treatment effect is null or negative.
We proxy cynicism using: (a) the pre-treatment item ``Adding people of color makes policing more fair'' (reverse-coded for cynical respondents) and (b) prior use-of-force contact.
We estimate heterogeneous effects within each of these subgroups.

\paragraph{Open-ended mechanism test.}
The survey includes an open-ended item (\texttt{Explain}) asking respondents to describe their reasoning.
We use a Bradley-Terry pairwise comparison model \citep{Bradley1952} to construct a latent trust score from GPT-assisted pairwise comparisons of responses, following the approach outlined in the main text.
We then estimate the treatment effect on the latent text-based trust score as a pre-registered robustness check.
Human-in-the-loop validation will be conducted on a random sample of 100 pairwise comparisons.

% ------------------------------------------------------------------
% D. MULTIPLE TESTING
% ------------------------------------------------------------------
\subsection{Multiple Testing and Pre-Specified Outcome Families}
\label{sec:multiple-testing}

We organize outcomes into three pre-specified families to control familywise error rates.

\begin{enumerate}
    \item \textbf{Primary family} (confirmatory): \texttt{Trust\_General}, \texttt{Trust\_Self}, \texttt{Trust\_Hiring}.
    We apply the Holm step-down correction \citep{Holm1979} within this family.
    All three outcomes must be corrected jointly for H1 and H2, yielding $m = 6$ tests total (3 outcomes $\times$ 2 hypotheses).
    We consider H1 (H2) supported if at least one corrected primary outcome is statistically significant in the expected direction, with the other outcomes showing consistent point estimates.

    \item \textbf{Secondary family} (pre-registered but not primary): \texttt{Procedural\_Justice}, \texttt{Police\_Resources}.
    We apply Holm correction within this family ($m = 4$ tests: 2 outcomes $\times$ 2 hypotheses).
    Significant results in this family are reported but do not constitute confirmatory evidence.

    \item \textbf{Exploratory} (no multiplicity correction): HTE / CATE results (Section~\ref{sec:hte}), mechanism tests (Section~\ref{sec:mechanisms}), open-ended text analysis.
    These results are clearly labeled exploratory in all tables and figures.
\end{enumerate}

% ------------------------------------------------------------------
% E. BALANCE CHECKS
% ------------------------------------------------------------------
\subsection{Randomization Validity and Balance Checks}

Prior to any outcome analysis, we will verify the integrity of block randomization by estimating an omnibus F-test of the null hypothesis that pre-treatment covariates jointly have no association with treatment assignment.
We regress each treatment indicator on the full set of pre-treatment covariates and report the $F$-statistic and $p$-value.
Individual covariate balance is reported in a standardized difference table (Table~1 in the main text), following the reporting conventions of \citet{Imbens2015}.
A $p$-value below 0.05 on the omnibus test will trigger a re-examination of the randomization procedure and a note in the paper; it will not alter the planned analysis.

We also report the distribution of treatment completion times across arms as a manipulation check for differential respondent engagement.

% ------------------------------------------------------------------
% F. ATTRITION AND EXCLUSION RULES
% ------------------------------------------------------------------
\subsection{Attrition, Exclusions, and Intent-to-Treat Estimand}

\paragraph{Primary estimand.}
The primary estimand is the \textit{intent-to-treat} (ITT) effect: the average difference in potential outcomes under treatment vs.\ control assignment, regardless of actual treatment uptake or engagement quality.

\paragraph{Pre-specified exclusion criteria.}
We will exclude respondents who:
\begin{enumerate}
    \item Complete the survey in fewer than 120 seconds (approximately one-third of the expected minimum completion time of $\sim$7 minutes), indicating inattentive responding;
    \item Fail any built-in attention check item (if included in the final instrument);
    \item Do not provide consent (screened out before randomization).
\end{enumerate}
These criteria are applied before outcome analysis.
All excluded respondents will be reported in a flow diagram following CONSORT guidelines.

\paragraph{Attrition analysis.}
If overall attrition exceeds 10\%, we will estimate differential attrition across arms using a logistic regression of completion on treatment assignment and covariates.
Differential attrition would threaten internal validity; we will report this prominently and, if detected, estimate Lee bounds \citep{Lee2009} as a sensitivity check.

\paragraph{Sensitivity analysis.}
All primary results will be re-estimated on the \textit{full} sample (no exclusions applied) as a robustness check, reported in the appendix.
